% =============================================================================
% §5 Model & sufficiency (M4 / R4) — GLASS BOX per Fina §5/§A:
% reframe (4-parameter cohort microsimulation), equations in main text,
% algorithm box, calibration in the open, pre-registered gates, wall,
% κ–φ separability, ablations, non-fit convergences (fig2 signature),
% M3 residual, division of labor restated. ODD appendix cross-ref.
% =============================================================================
\section{A Transparent Sufficiency Test: Is Cohort Replacement Enough?}
\label{sec:model}

Section~\ref{sec:discrimination} identified the mechanism. This section asks a
different question, and it is important to keep the two separate: taking the
identified structure as given, is it \emph{sufficient}? Can cohort replacement
alone --- with no reflexivity anywhere --- generate the observed composition
collapse, its magnitude \emph{and} its late-loaded shape? The simulation model
below is never asked to identify anything; identification is econometric and
lives in Section~\ref{sec:discrimination}. The model's job is to test whether
the surviving mechanism quantitatively suffices, under an evaluation whose
success criteria were fixed before the model was run.

\subsection{The model is four interpretable parameters, not a black box}
\label{sec:modelspec}

The model is best described as a parsimonious cohort microsimulation. Beyond
five baseline demographic rates (calibrated, and interpretable as ordinary
formation/dissolution hazards), its entire mechanism content is \emph{four}
parameters: three describing how marriage propensity declines across birth
cohorts, and one describing how union formation tilts across birth cohorts.
There are no per-cohort free effects, no adaptive dynamics, no interaction
terms, and --- by construction --- no channel through which any agent's
transition depends on any other agent's state.

Each woman $i$ carries age $a_{it} \in \{15,\dots,49\}$, birth year $b_i$, and
union status $u_{it} \in \{\Sng, \Coh, \Mar\}$. One tick is one calendar year.
Annual transitions are:
\begin{align}
\Pr(\Sng \to \text{union})_{it} &= f_0 \,\phib[b_i],
  &&\text{of which married w.p. } m_0\,\kb[b_i], \text{ else cohabiting},
  \label{eq:formation}\\
\Pr(\Coh \to \Sng)_{it} &= \delta_c, \qquad
\Pr(\Coh \to \Mar \mid \text{no dissolution})_{it} = c_0 \,\kb[b_i],
  \label{eq:conversion}\\
\Pr(\Mar \to \Sng)_{it} &= \delta_m.
  \label{eq:dissolution}
\end{align}
In the cohabiting state, where dissolution and conversion compete, the two
risks are evaluated \emph{sequentially} --- dissolution first, then conversion
conditional on remaining in union, exactly as in Algorithm~\ref{alg:abm} ---
so the realized annual conversion probability is
$(1-\delta_c)\,c_0\,\kb$. This matches the reference implementation; the
conversion parameter $c_0$ is interpreted conditionally throughout.
The two cohort terms are
\begin{align}
\kb &= \kappa_\infty + \frac{1-\kappa_\infty}{1+\exp\{s\,(b-b_0)\}},
  \label{eq:kappa}\\
\phib &= \exp\{-\tau\,(b-1985)\}, \quad \text{clamped to } [0.2,\,3].
  \label{eq:phi}
\end{align}
Equation~\eqref{eq:kappa} is the paper's mechanism: a logistic decline in
cohort marriage propensity with exactly three parameters --- a floor
$\kappa_\infty$ (the marriage propensity of the newest cohorts relative to the
oldest), a midpoint $b_0$ (the birth cohort at which the decline is halfway),
and a slope $s$ (how compressed the across-cohort change is). It multiplies
both marriage-entry margins --- the formation split in~\eqref{eq:formation} and
the conversion hazard in~\eqref{eq:conversion} --- and nothing else.
Equation~\eqref{eq:phi} is a one-parameter log-linear tilt in overall union
formation across cohorts. Old cohorts ($b \ll b_0$) have $\kb \approx 1$: the
baseline rates $\{f_0, m_0, c_0, \delta_c, \delta_m\}$ are simply the behavior
of pre-decline cohorts.

In a robustness specification we add an exogenous period shifter
$\pit = 1 - d\,/\,(1+\exp\{-s_p (t - t_p)\})$ multiplying the same
marriage-entry margins --- calendar time only, common to all cohorts, never a
function of prevailing composition (the Colombian period overlay of
Section~\ref{sec:discrimination}). As shown below, the data render it
dispensable, and the headline specification omits it entirely.

Algorithm~\ref{alg:abm} states the per-year update completely; the ODD protocol
appendix (Appendix~\ref{app:odd}) documents the model to the community
transparency standard \citep{grimm2020odd}, sufficient for reimplementation.

\begin{algorithm}[t]
\caption{One simulated year (population of $N$ women; one tick = one year)}
\label{alg:abm}
\begin{algorithmic}[1]
\For{each woman $i$ (order immaterial; no interaction)}
  \State $\kappa \gets \kb[b_i]$;\quad $\varphi \gets \phib[b_i]$
    \Comment{depends only on own birth year}
  \If{$u_i = \Sng$}
    \State with prob.\ $f_0\varphi$: form union; married w.p.\
      $m_0\kappa$, else cohabiting
  \ElsIf{$u_i = \Coh$}
    \State with prob.\ $\delta_c$: $u_i \gets \Sng$; else with prob.\
      $c_0\kappa$: $u_i \gets \Mar$
  \Else
    \State with prob.\ $\delta_m$: $u_i \gets \Sng$
  \EndIf
  \State $a_i \gets a_i + 1$
  \If{$a_i > 49$} \State replace: $a_i \gets 15$, $u_i \gets \Sng$,
    $b_i \gets t - 15$ \Comment{new entry cohort carries its own $\kb$}
  \EndIf
\EndFor
\State $t \gets t + 1$ \Comment{model step advances the calendar; nothing else}
\end{algorithmic}
\end{algorithm}

The no-reflexivity claim is verifiable in the code, not just the prose: the
agent update reads no other agent's state (no population shares, no
reference-group terms enter any transition), and the model-level step advances
the calendar only. Automated checks --- static source checks plus a runtime
registry of every data file opened during calibration --- are logged with the
repository (eleven checks per country, all passing).

\subsection{Calibration in the open}
\label{sec:calibration}

\emph{Targets.} The observed per-band (20--24 through 35--39) married and
cohabiting shares, all available years: Costa Rica ENAHO 2010--2024; Colombia
GEIH 2008--2024 (2008 seed per the documented 2007 frame outlier,
Section~\ref{sec:datalim-measurement}). The initial population is seeded at
the first observed year's composition.

\emph{Loss.} Composition only:
\begin{equation}
\mathcal{L} \;=\; \sum_{\text{bands}} \sum_{\text{years}}
\left[ \left(m^{\text{sim}} - m^{\text{obs}}\right)^2
     + \left(c^{\text{sim}} - c^{\text{obs}}\right)^2 \right].
\label{eq:loss}
\end{equation}

\emph{Identification wall.} The TFR never enters the loss, tunes no parameter,
and is not loaded during calibration (enforced by the runtime file registry).
Where a TFR path appears in this paper it is a comparison overlay only.

\emph{Free versus fixed.} Free, per country: the five baseline rates, the three
$\kb$ parameters, the $\phib$ tilt (and, in the robustness specification, the
three $\pit$ parameters). Fixed: the demographic scaffolding (age range,
seeding, closed-cohort replacement), reported in Appendix~\ref{app:odd}. The
resulting parameter-to-moment ratio is the most legible overfitting check:
\emph{nine} free parameters per country face 120 (Costa Rica: 4 bands
$\times$ 2 shares $\times$ 15 years) and 136 (Colombia: $4 \times 2 \times 17$)
target moments. Closed-cohort replacement holds cohort inflow constant ---
counterfactual under birth declines of 18--47 percent --- but the evaluation
metrics are computed on \emph{within-band composition shares}, which are
insulated from inflow scale by construction.

\emph{Optimizer.} Nelder--Mead in a boxed transform, four random restarts, 600
evaluations per country; inner evaluation ensembles of 12{,}000 agents
$\times$ 4 seeds; final reported ensembles of 50{,}000 agents $\times$ 16 seeds
(mean $\pm$ s.d.). The pipeline is deterministic given seeds; the Costa Rican
calibration, re-run end-to-end, reproduces its loss to all reported digits.

\emph{Pre-registered evaluation.} Three metrics were fixed, with their
thresholds, before the full calibration ran. On the equal-band-mean married
share: \textbf{M1}, magnitude capture --- the simulated share of the observed
seed-to-2024 fall --- must be at least $0.85$; \textbf{M2}, the late-window
share of the total decline (2017--2024 relative to the full window) must match
the observed within $\pm 0.10$; \textbf{M3}, the sign of the mean second
difference over 2015--2024 must match (does the decline accelerate or
decelerate late?).

\subsection{Results: sufficiency on magnitude and late-loading}
\label{sec:results}

\begin{table}[t]
\centering
\caption{Pre-registered sufficiency metrics, headline (cohort-only)
specification. Final ensembles, 50{,}000 agents $\times$ 16 seeds.}
\label{tab:gates}
\begin{tabular}{lccl}
\toprule
Metric (married, 20--39 equal-band) & Costa Rica & Colombia & Criterion \\
\midrule
M1 magnitude capture & 0.954 & 0.999 & $\geq 0.85$: \textbf{pass, pass} \\
M2 late-share, sim vs.\ obs & 0.490 / 0.581 & 0.468 / 0.459 & $|\Delta|\leq0.10$: \textbf{pass, pass} \\
M3 late-window $\Delta^2$ sign & $+$ / $+$ & $+$ / $-$ & same sign: pass, \textbf{fail} \\
\midrule
Composition loss~\eqref{eq:loss} & 0.0372 & 0.0246 &
  (threshold model: 0.0855) \\
\bottomrule
\end{tabular}
\end{table}

Table~\ref{tab:gates} reports the gates; Figure~\ref{fig:abmbands} the fit.
The verdict, stated precisely: \emph{sufficiency on magnitude and late-loading,
with one unreproduced shape residual.} The cohort-only model captures 95.4
(Costa Rica) and 99.9 (Colombia) percent of the observed married-share fall and
matches the late-loading of the decline within 0.09 and 0.01 respectively. The
residual is Colombia's M3: the observed series steepens slightly in the late
window (a 2019--22 kink in which both married and cohabiting shares break
trend) while the simulated path does not. Both second differences are near
zero; the kink coincides with pandemic-era survey collection
(Section~\ref{sec:datalim-measurement}) and/or a discrete shock, and its
resolution requires transition-level data
(Section~\ref{sec:limitations}). The miss is structural, not incidental: both
cohort terms are monotone in birth cohort --- $\phib$ in particular cannot
produce Colombia's flat-then-declining total-union path --- so the same
parsimony that disciplines the model is what surfaces here as the residual. A kink that a smooth model misses is not
evidence of a cascade; Section~\ref{sec:discrimination} tested for that
signature directly and rejected it. Honest margins: under the headline
specification Costa Rica's M2 sits nearer its tolerance than with the period
term included ($\Delta$ 0.091 versus 0.038), and the (ungated) cohabitation
magnitude capture drops to 0.63.

\begin{figure}[t]
\centering
\begin{subfigure}[b]{\textwidth}
  \includegraphics[width=\textwidth]{fig_abm_bands_COL.pdf}
  \caption{Colombia (seed 2008)}
\end{subfigure}\\[0.6em]
\begin{subfigure}[b]{\textwidth}
  \includegraphics[width=\textwidth]{fig_abm_bands_CRI.pdf}
  \caption{Costa Rica (seed 2010)}
\end{subfigure}
\caption{Headline (cohort-only) model versus observed union composition, by
five-year age band: married (top rows) and cohabiting (bottom rows) shares,
simulated mean $\pm 1$ s.d.\ over 16 seeds against observed.}
\label{fig:abmbands}
\end{figure}

\subsection{Why the result is credible: what a black box could not do}
\label{sec:credibility}

Fit quality is table stakes --- a model handed a flexible mechanism should fit
a trajectory. The evidential weight lies in three properties a curve-fitting
box would not exhibit.

\paragraph{(i) A single recovered parameter matches an estimate the model never
saw.} The calibrated $\kb$ midpoints land at birth cohorts \textbf{1991.7}
(Costa Rica) and \textbf{1988.8} (Colombia). The APC decomposition of
Section~\ref{sec:discrimination} --- estimated by entirely different methods,
on cohort effects the simulation never touched --- places the married cohort
effects' sign crossover in the same 1985--1992 window
(Figure~\ref{fig:kappaapc}). Because $\kb$ is a three-parameter logistic, the
midpoint is one number, not a flexible profile that could have been bent to
agree: the convergence is non-circular. This is the paper's signature exhibit.

\begin{figure}[t]
\centering
\begin{subfigure}[b]{0.48\textwidth}
  \includegraphics[width=\textwidth]{fig_kappa_apc_COL.pdf}
  \caption{Colombia}
\end{subfigure}
\hfill
\begin{subfigure}[b]{0.48\textwidth}
  \includegraphics[width=\textwidth]{fig_kappa_apc_CRI.pdf}
  \caption{Costa Rica}
\end{subfigure}
\caption{The signature exhibit. Top panels: the calibrated cohort multiplier
$\kb$ (three parameters, calibrated to composition trajectories only). Bottom
panels: the independently estimated APC married cohort effects
(Section~\ref{sec:discrimination}). The calibrated midpoint recovers the APC
sign crossover the model was never fit to.}
\label{fig:kappaapc}
\end{figure}

\paragraph{(ii) Ablations reproduce the econometric decomposition.} Removing a
mechanism and \emph{fully recalibrating} the rest yields losses that mirror
Section~\ref{sec:discrimination}: without the cohort gradient, Colombia's loss
rises 63 percent (cohort replacement is essential exactly where the APC found
it dominant); without the period shifter, the loss is essentially unchanged in
both countries (Costa Rica $+6$ percent; Colombia slightly \emph{better}).
The apparent tension with Colombia's statistically significant period effect
dissolves on inspection: that effect is real but small (Section
\ref{sec:verdict}), and statistical existence is not quantitative necessity.
This is why the cohort-only model is the headline specification and the period
shifter a documented robustness term. A corollary worth stating against
ourselves: since a period-only model can also pass the magnitude gates,
band-level composition alone cannot discriminate the mechanisms --- which is
precisely why identification lives in Section~\ref{sec:discrimination} and the
model is asked only about sufficiency.

\paragraph{(iii) The mechanism's parameters are pinned by separate moments.}
Both cohort terms act on union behavior, so one may ask whether they are
separately identified or a disguised degree-of-freedom. They are separable by
construction: $\kb$ multiplies only marriage-\emph{entry} margins, both of
which leave union status unchanged (a formation splits into married or
cohabiting either way; a conversion moves within union), so $\kb$ is
union-total-neutral; $\phib$ multiplies formation itself and moves the union
total directly. Paired-seed perturbations at the calibrated parameters confirm
it: perturbing $\kappa_\infty$ moves the married share with the union total
essentially unchanged (response ratios 11.6 in Costa Rica, 102 in Colombia),
while perturbing $\tau$ does the reverse (ratios 0.32 and 0.05). The married
split pins $\kappa$; the union total pins $\varphi$.

Finally, the comparison that frames the exercise: this four-parameter,
no-feedback model fits the composition data \emph{better} than the calibrated
coordination-threshold model it replaced (losses 0.037/0.025 versus 0.086 on
the identical target) --- the reflexive mechanism was not removed at the cost
of fit, but with a gain.

A comparison-only overlay of the generated fertility path against the national
TFR appears in the online appendix figures; consistent with the identification
wall, we attach no claim to it --- its level rides a placeholder fertility
schedule, and the fertility-magnitude question requires union-specific
fertility rates that are the subject of ongoing data acquisition
(Section~\ref{sec:limitations}).
